% RESULTADO_RECUPERADO from II REV09. See MAPA_PROCEDENCIA.json.
% Selected lines and explicit mechanical changes; no new proof.
\section{Hexadic--octadic incidence and the Barbero--Immirzi functional}
\label{v:ant:sec:barbero}
Let $H\hookrightarrow O$ be a marked inclusion, $|H|=6$, $|O|=8$.
The incidence objects
\[
\mathcal F_6=H\times\binom H2,\qquad
\mathcal F_8=O\times\binom H2
\]
have cardinalities $90$ and $120$. Their normalized difference gives
\[
\frac{|\mathcal F_6|-|\mathcal F_8|}{24\binom62}=-\frac1{12}.
\]
\begin{figure}[htbp]
\centering
\begin{tikzpicture}[>=Latex,every node/.style={font=\small}]
\draw[rounded corners,draw=ink] (-1.8,-1.35) rectangle (1.8,1.35);
\foreach \a in {0,60,...,300}{
\fill[ink] (\a:0.83) circle(2pt);}
\node at (0,-1.7) {$H:\ 6$ points};
\draw[->,thick,ink] (2.1,0)--(3.6,0);
\node[above] at (2.85,1.5) {marked inclusion};
\draw[rounded corners,draw=ink] (4,-1.35) rectangle (7.6,1.35);
\foreach \a in {0,60,...,300}{
\fill[ink] ($(5.8,0)+(\a:0.83)$) circle(2pt);}
\fill[accent] (4.5,0) circle(2pt);
\fill[accent] (7.1,0) circle(2pt);
\node at (5.8,-1.7) {$O:\ 8$ points};
\node at (0,-2.35) {$6\binom62=90$};
\node at (5.8,-2.35) {$8\binom62=120$};
\end{tikzpicture}
\caption{Marked incidence from six to eight points. The drawing represents
the sets and the inclusion; it does not purport to turn the faces or vertices
of a cube into Steiner blocks without their incidence map.}
\end{figure}

\subsection{The weights precede the pole coefficient}
In the free module of oriented events, the source fixes the
fundamental chain of the turn
\[
c_{12}^{+}=
\sum_{j\in\Z/12\Z}[j,\mathcal F_6]
-[\partial_{\rm ret},\mathcal F_8].
\]
Multiplicity twelve corresponds to the sectors, and multiplicity one
to the oriented seam. The involution changes the orientation of the complete
chain. Let
\[
S_m(q)=\frac{q^m}{1-q^{3m}},\qquad 0<q<1.
\]
The linear reader assigns $S_{90}$ to each sector and $S_{120}$ to the seam;
consequently,
\[
\mathscr R(c_{12}^{+})=12S_{90}-S_{120}.
\]

\begin{theorem}[Pole coefficient and conjugate evaluation]
The image of the fundamental chain has coefficient $1/24$
with respect to $(1-q)^{-1}$ and residue $-1/24$ with respect to $(q-1)^{-1}$.
Its evaluation in the previously generated channels determines
\begin{equation}
\gamma_{\HMT}=
\frac{\pi AC^*}{180}
+12\bigl[S_{90}(q_-)-S_{90}(q_+)\bigr]
-\bigl[S_{120}(q_-)-S_{120}(q_+)\bigr].
\label{v:ant:eq:gamma}
\end{equation}
\end{theorem}
\begin{proof}
The factorization $1-q^{3m}=(1-q)\sum_{j=0}^{3m-1}q^j$ implies
$\lim_{q\uparrow1}(1-q)S_m(q)=1/(3m)$. By linearity,
\[
\frac{12}{270}-\frac1{360}=\frac1{24}.
\]
The coordinate $q-1$ reverses the sign of the residue.
Conjugate evaluation and the angular component of the source
produce \eqref{v:ant:eq:gamma}.
\end{proof}

The Diophantine certificate $4a-3b=45$ admits the pairs
$(12+3t,1+4t)$, $t\ge0$; it certifies minimality of $(12,1)$,
but does not replace the origin of the multiplicities in the chain.
The evaluation of \eqref{v:ant:eq:gamma} is
\[
\gamma_{\HMT}
=0.27400767269445927474725526077398041940\ldots.
\]
This number is an output of the source functional; it is not an external value
used to select $A$, $C^*$, or their coefficients.

\subsection{Parity and structural sensitivity}
Reversing $C^*$ exchanges $q_+$ and $q_-$ and changes the sign
of the angular component. Hence
\[
\gamma_{\HMT}(A,-C^*)=-\gamma_{\HMT}(A,C^*).
\]
The functional preserves orientation, whereas other readers,
such as $q_+q_-$, are even. This difference precedes their physical interpretation.
